\documentclass[11pt]{article}
\usepackage{fontspec}
\setmainfont{Latin Modern Roman}
\usepackage{amsmath,amssymb,mathtools,unicode-math}
\setmathfont{Latin Modern Math}
\usepackage{xurl,hyperref}
\hypersetup{hidelinks,pdftitle={Problem statement}}
\newcommand{\sref}[2]{\hyperlink{source:#1:#2}{[S#2]}}
\newcommand{\cataloguescope}[1]{\paragraph{Mathematical question.}#1}
\begin{document}
\section{Cores of strongly regular graphs}\label{statement}

\subsection{Definitions and mathematical statement}

\paragraph{Shared notation.}$\mathbb N=\{1,2,\ldots\}$, and $\mathbb Z,\mathbb Q,\mathbb R,\mathbb C$ denote the integers, rationals, reals and complex numbers. Any other conventions are specified locally.


A homomorphism $f:G\to H$ of finite simple graphs is a vertex map satisfying $f(u)f(v)\in E(H)$ whenever $uv\in E(G)$. A core of $G$ is a graph of minimum vertex cardinality among graphs admitting homomorphisms both from and to $G$; it is unique up to isomorphism. A graph is strongly regular with parameters $(v,k,\lambda,\mu)$ if it has $v$ vertices, every vertex has degree $k$, each adjacent pair has $\lambda$ common neighbors and each distinct nonadjacent pair has $\mu$ common neighbors. A complete graph $K_s$ joins every distinct vertex pair.
\paragraph{Statement recorded in the supplied source.}
Does every strongly regular graph have a core isomorphic either to itself or to a complete graph? \sref{statement}{2}
\subsection{Short English statement}
Is the core of every strongly regular graph either the whole graph or a clique?
\subsection{Sources}
\begin{description}
\item[\hypertarget{source:statement:1}{S1}] ULAM.AI and the UnsolvedMath Contributors, UnsolvedMath: A Curated Collection of Open Mathematics Problems, catalogue entry. CC BY 4.0. \url{https://huggingface.co/datasets/ulamai/UnsolvedMath}.
\item[\hypertarget{source:statement:2}{S2}] Open Problem Garden, Cores of strongly regular graphs, original node 824, statement and mathematical discussion. \url{https://www.openproblemgarden.org/op/cores_of_strongly_regular_graphs}.
\end{description}
\label{end:statement}
\end{document}
